\section{Conclusion}
\label{sec:conclusion}

We presented the ANU team's submission to PlantCLEF 2026. Our best public Macro F1 of \textbf{0.418265} and private Macro F1 of \textbf{0.40283} were obtained from a BioCLIP 2.5 ViT-H/14 fine tune with only the last four transformer blocks (together with the final layer norm and projection) unfrozen, trained for ten epochs on the redownloaded PlantCLEF 2024 manifest (2.65 million single-plant images across 7,806 species, with pre-filled genus and family labels) under a joint species/genus/family cross entropy loss. Inference combined a four-by-four tile decomposition at tile size 448 with no overlap, softmax mean aggregation across the sixteen tiles, an ensemble of inference at 224 and 336 pixel resolutions, logit adjustment ($\tau{=}0.25$), and an adaptive probability threshold of $T{=}0.03$. A follow-up experiment (i003) added 161,686 extra images for the 2,022 species with fewer than one hundred training examples and capped every species at five hundred training images; the resulting checkpoint reached 0.40041 on the public leaderboard, below the uncapped anchor by 0.018, which we read as evidence that on this benchmark the marginal value of additional head-class examples is positive and that flattening the head discards information rather than rebalancing it.

Three findings carry the bulk of what this benchmark taught us. The first is that partial unfreeze is a sharp, not a plateau, design choice on BioCLIP 2.5: $n{=}4$ unfrozen blocks is the empirical sweet spot, deviations of one block in either direction cost between $0.009$ and $0.014$ Macro F1, and a full fine-tune destroys the Tree-of-Life prior and drops the model to 0.208. The second is the validation versus Kaggle decoupling: across our twelve strongest runs the two metrics correlate at only $r \approx 0.4$, and within the unfreeze ablation they invert entirely, which is direct evidence that the bottleneck on this benchmark is the single-plant to multi-species distribution shift rather than model capacity. The third is that auxiliary post-hoc reweighting of an already strong visual posterior carries essentially no usable signal beyond what the visual model has already absorbed. Classifier retraining over a frozen triple-backbone fusion (cRT, $\Delta {=} -0.036$) introduced wrong evidence and hurt; genus and family co-occurrence priors derived from the same posterior ($\Delta {=} -0.006$) and a seasonal phenology prior over GBIF observation dates ($\Delta {=} -0.005$) both moved the anchor by amounts close to our run-to-run noise floor. Fine-tuned visual classifiers on biological data absorb these signals implicitly through training-distribution co-occurrence, leaving little for post-hoc combination to exploit.

Several directions for future work follow directly from these findings. The first is to pull more aggressively on the data axis rather than the capacity axis: a synthetic mosaic training regime that composes single-plant images into multi-species scenes with controlled species overlap would let the network see the quadrat distribution during training, which neither the PC24 corpus nor the iNaturalist habitat shots do at scale. The second is self-distillation against our own production checkpoint, using its predictions on the unlabelled LUCAS pseudo-quadrat corpus as soft targets filtered by confidence and entropy. The third is genuinely orthogonal evidence that is not in the image at all, in particular a location-conditioned bioclimatic envelope prior built from training-image coordinates against WorldClim variables at each test site --- the one orthogonal axis that visual features cannot recover from the image alone. Finally, per-class threshold calibration on a held-out multi-species validation set would more directly target the Macro F1 metric than the global probability threshold we currently use, and could absorb residual species-frequency imbalances that our manifest cap leaves behind.
